% §3.1 Suite and Kernel Selection
%   - First: scope of the survey — 35 HPC repos, multiple categories.
%   - Then: KEY INSIGHT — repo-level counting understates benchmark
%     material by 79×. Kernel-level counting reveals 472 CUDA–OMP pairs
%     from just 6 repos. This motivates kernel-centric evaluation.
%   - Then: selection criteria (stated briefly) — multi-API equivalence,
%     build-run-verify automation, self-checking output, domain diversity.
%   - Then: describe each selected suite in order of corpus contribution:
%     * Rodinia (22 kernels, 60 specs — largest contributor)
%     * HeCBench (10 curated from 516 via structured funnel)
%     * XSBench, RSBench, mixbench (smaller, domain-diverse)
%
% §3.2 Evaluation Corpus
%   - First: combined table showing suite composition + complexity metrics.
%   - Then: headline statistics — 87 verified-PASS specs, SLoC range
%     (80–3,304), 89% exceed ParEval-Repo's 133 SLoC threshold.
%   - Finally: multi-file fraction (25%) — sets up the finding in Results
%     that file-boundary coordination is an independent difficulty axis.

\section{Benchmark Curation}
\label{sec:benchmark-curation}

The benchmark corpus was assembled through a four-stage systematic process: surveying open-source HPC benchmark repositories, quantifying kernel-level translation opportunities across parallel APIs, filtering candidate kernels through automated build-run-verify checks, and verifying each selected kernel through the complete pipeline on the evaluation platform. This curation populates the Source Code component of the ParBench architecture (Figure~\ref{fig:architecture}, orange) with kernel files from five benchmark suites that collectively span the CUDA--OpenMP--OpenCL API triple.

\subsection{Suite Selection}
\label{sec:suite-selection}

\textbf{Survey methodology.} We conducted a systematic survey of open-source HPC benchmark repositories, beginning with 40 candidate archives identified from the ECP proxy application catalog, published benchmark suites, and HPC conference proceedings. After excluding 5 repositories for download failures or insufficient documentation, 35 repositories spanning six benchmark types: suites, mini-applications, proxy applications, full applications, libraries, and microbenmarks were analyzed in detail. Some examples include HeCBench~\cite{HeCBench2023}, Rodinia~\cite{Rodinia2009}, BabelStream, CloverLeaf, SW4lite, LAMMPS, Kokkos Kernels; the full repository list is in Appendix~\ref{sec:appendix-b}.
%Of the 35, 30 were classified Tier~A (documented build, automated verification, active maintenance) and 5 as Tier~B (partial verification or limited API coverage).

\textbf{Repository-level vs.\ kernel-level counting.} A central finding of the survey is that repository-level counting dramatically understates the available benchmark material for translation evaluation. The API co-occurrence matrix identifies 6 repositories containing both CUDA and OpenMP implementations. However, kernel-level analysis, enumerating individual computational kernels that have verified equivalent implementations across APIs, reveals 472 independent CUDA--OpenMP translation pairs across those same 6 repositories, a 79$\times$ multiplier (Figure~\ref{fig:repo-vs-kernel}). This extreme concentration is driven by large multi-kernel suites: HeCBench alone contributes 325 CUDA--OpenMP kernel pairs, RAJAPerf contributes 106, and Rodinia contributes 19. The same pattern holds across other API pairs: 633 CUDA--HIP kernel pairs (across 3 repositories, a 211$\times$ multiplier) and 616 CUDA--SYCL pairs (across 2 repositories, a 308$\times$ multiplier). This finding motivates a kernel-centric evaluation strategy in which benchmarks are evaluated at the granularity of individual computational kernels rather than entire repositories.
% [Float moved to Appendix D]
%The repository-to-kernel extraction ratio is shown in Appendix Figure~\ref{fig:repo-vs-kernel}.

% [Float moved to Appendix D]
%Translation pair counts per kernel are detailed in Appendix Table~\ref{tab:kernel-pairs}.

\textbf{Selection criteria.} Five criteria guided suite selection: (1)~\textbf{open-source availability} for reproducibility;(2)~\textbf{multi-API kernel equivalence}, requiring implementations of the \emph{same} kernel across multiple APIs to ensure authoritative reference translations rather than independently developed programs; (3)~\textbf{build-run-verify automation} each kernel must be buildable, runnable, and verifiable without human intervention; (4) \textbf{Self-checking correctness.} Kernels must produce self-checking output--deterministic checksums, tolerance-bounded numerical comparison, or labeled correctness indicators (e.g., \texttt{PASS}/\texttt{FAIL}, \texttt{verify()}) - enabling automated correctness verification without external reference files. (5)~\textbf{domain diversity} across computational patterns. Applying them to the 35 surveyed repositories yielded five suites that satisfy all five criteria and therefore collectively provide maximum coverage of the CUDA--OpenMP--OpenCL API triple. %These criteria are intentionally conservative: they exclude repositories that require interactive execution, external datasets not bundled with the source, or proprietary licenses.


\subsection{Kernel Selection}
\label{sec:kernel-selection}

Kernel selection proceeds in two stages: selection from HeCBench (which provides the largest pool of multi-API kernels) and curation of kernels from Rodinia and three additional suites.

\textbf{HeCBench selection funnel.} HeCBench~\cite{HeCBench2023}\footnote{Commit \texttt{22785cdd}.} is the largest source of multi-API kernel implementations in the survey, with 516 kernels spanning CUDA, HIP, SYCL, and OpenMP. We implemented a 5-stage structured selection funnel to identify kernels suitable for automated translation evaluation. The funnel filterered for \textbf{4-API completeness} (327 of 516 kernels provide implementations in CUDA, HIP, SYCL, OpenMP APIs); \textbf{Build system presence} (325 of these 327 have Makefiles present in the CUDA variant);\textbf{Self-checking output patterns} (242 of the 325 contain self-checking output patterns for string matching-based automated correctness verification in the source for \texttt{PASS}, \texttt{FAIL}, \texttt{verify}, or \texttt{correct});\textbf{Complexity bounds} (excluding kernels with  $>$15 or $<$2 source files or kernels requiring external input data files, unbundled to source); and, \textbf{Domain diversity} (60 kernels were selected to maximize coverage across computational domains).

From this 60-kernel working set, a curated subset of \textbf{10 kernels} was selected as those with verified OMP-target GPU offload variants in addition to CUDA, maximizing domain diversity: stencil computation (stencil1d, heat2d, iso2dfd), graph algorithms (floydwarshall, page-rank), combinatorial search (nqueen), molecular dynamics (md), signal processing (convolution1d, scan), and numerical methods (jacobi). The curated subset provides 25 specs across three API variants: all 10 have CUDA and OMP-target GPU offload implementations, and 5 additionally have CPU OpenMP variants.

\textbf{Rodinia.} Rodinia~\cite{Rodinia2009}\footnote{Commit \texttt{9c10d3ea}.} is the largest contributor to the evaluation corpus, providing CUDA, OpenMP, and OpenCL implementations of most kernels in the same source tree, yielding authoritative reference implementations across the full API triple. ParBench includes 60 Rodinia specs covering 22 kernels (22 CUDA, 18 OpenMP, 20 OpenCL; non-uniform because not all kernels have all three API variants). After systematic verification, 53/60 achieve PASS and 7 are \knownfail{} for platform-specific build, runtim eor verification failures. However, Rodinia's wide adoption raises a data contamination concern; the AST-driven augmentation engine (Section~\ref{sec:augmentation-engine}) mitigates this by generating structurally novel variants that preserve semantics while defeating memorization.

% \begin{itemize}
% \item 2 specs fail because CUDA~12 removed the deprecated \texttt{texture<>} reference API (kmeans-cuda, mummergpu-cuda);
% \item 1 spec fails due to a missing OpenGL dependency (hybridsort-cuda);
% \item 1 spec inherits CUDA texture dependencies in its OpenMP variant (mummergpu-omp);
% \item 2 OpenCL specs exhibit pre-existing runtime crashes (nn-opencl, kmeans-opencl).
% \end{itemize}

% These failures are documented but excluded from the evaluation corpus; they reflect toolchain evolution, not benchmark design defects.

%However, Rodinia's age and wide availability raise a legitimate concern: its source code is likely present in LLM training corpora, potentially inflating pass rates through memorization rather than genuine translation competence. To address this threat to validity and broaden domain coverage, ParBench includes three additional benchmark suites.

\textbf{XSBench}~\cite{XSBench2014}\footnote{Commit \texttt{ba08e522}.} (4 specs: CUDA, OpenMP, OpenCL, OMP-target; all PASS) is a Monte Carlo neutron transport proxy application that implements the continuous-energy macroscopic cross-section lookup kernel from OpenMC. XSBench is also evaluated in prior repository-level work~\cite{ParEvalRepo2025}, enabling direct comparison of evaluation granularity on the same kernel.

\textbf{RSBench and mixbench.} RSBench~\cite{RSBench2014}\footnote{Commit \texttt{34b64478}.} (4 specs: CUDA, OpenMP, OpenCL, OMP-target; all PASS) is a reactor simulation proxy dervide from the same OpenMC cross-section lookup as XSBench but using a multi-pole method with complex arithmetic and Faddeeva function evaluation. mixbench~\cite{mixbench2017}\footnote{Commit \texttt{32edeca9}.} (3 specs:CUDA, OpenMP, OpenCL; all PASS) is a GPU micro-benchmark measuring compute-to-memory-bandwidth balance. Together they contribute domain diversity beyond Rodinia's traditional HPC kernels.

% \subsection{API Coverage}
% \label{sec:api-coverage}

% \begin{table*}[!htbp]
% \centering
% \caption{Survey---kernel-level translation pair counts.}
% \label{tab:kernel-pairs}
% \small
% \begin{tabular}{@{}lccl@{}}
% \toprule
% API Pair & Repos & Kernel Pairs & Primary Sources \\
% \midrule
% CUDA--HIP & 3 & 633 & HeCBench (504), RAJAPerf (106), CloverLeaf (16) \\
% CUDA--SYCL & 2 & 616 & HeCBench (487), RAJAPerf (106), CloverLeaf (16) \\
% CUDA--OpenMP & 6 & 472 & HeCBench (325), RAJAPerf (106), Rodinia (19) \\
% HIP--OpenMP & 2 & 453 & HeCBench (324), RAJAPerf (106), CloverLeaf (16) \\
% CUDA--OpenCL & 6 & ${\sim}$200 & Rodinia (19), Parboil, SHOC \\
% \bottomrule
% \end{tabular}
% \end{table*}

% The survey data directly inform ParBench's API selection. CUDA serves as the primary source language, reflecting its dominant position in GPU programming. OpenMP is the primary translation target: the kernel-level survey identified CUDA-to-OpenMP as the largest translation opportunity among CPU-targeting APIs, with 472 kernel pairs across 6 repositories (Table~\ref{tab:kernel-pairs}).

% Although CUDA--HIP and CUDA--SYCL yield higher raw kernel-pair counts (633 and 616 respectively; Table~\ref{tab:kernel-pairs}), these pairs represent a qualitatively different---and substantially easier---translation challenge. HIP is designed as a near-syntactic mirror of CUDA: \texttt{cudaMalloc} maps to \texttt{hipMalloc}, kernel launch syntax is preserved verbatim, and thread-index arithmetic transfers unchanged. In contrast, CUDA-to-OpenMP translation requires structural program transformation: SPMD thread-index arithmetic must be refactored into fork-join loop parallelism, explicit GPU memory management must be eliminated entirely, and synchronization primitives (\texttt{\_\_syncthreads()}) must be replaced with implicit OpenMP barrier semantics or explicit \texttt{critical}/\texttt{atomic} sections. This paradigm gap makes CUDA--OpenMP the most scientifically informative translation direction for evaluating whether LLMs reason about parallel structure rather than performing surface-level API renaming.
% % [Float moved to Appendix D]
% API characteristics and execution model comparison are detailed in Appendix Table~\ref{tab:api-characteristics}.


% OpenCL provides a secondary translation target that exercises a qualitatively different programming model from both CUDA and OpenMP. Rodinia's particular strength lies in its OpenCL coverage: 20 of 22 Rodinia kernels have OpenCL implementations, compared to only sparse OpenCL coverage in HeCBench and RAJAPerf. This makes Rodinia the primary source for OpenCL translation pairs.

% OpenMP target offload (OMP-target) provides a fourth API that uses compiler-directed GPU offloading via \texttt{\#pragma omp target}. It is available for XSBench, RSBench, and the HeCBench curated kernels, and requires the NVIDIA HPC compiler (\texttt{nvc}) rather than standard GCC. Because \texttt{nvc} is not universally available and OMP-target's compilation model differs substantially from CPU OpenMP, OMP-target directions are evaluated as case studies rather than as part of the standard evaluation.
\subsection{Evaluation Corpus}
\label{sec:eval-corpus}

\begin{table*}[!htbp]
\centering
\caption{Suite summary showing total specs, verified PASS, KNOWN\_FAIL (KF), and API coverage for each of the five benchmark suites.}
\label{tab:suite-summary}
\small
% src: Verified 2026-04-03 against specs/ on disk:
% src: Rodinia: ls specs/rodinia-*.json | wc -l => 60; unique kernels => 22; KF=7 ([known-issues.md](http://known-issues.md))
% src: HeCBench: 10 curated kernels x {cuda,omp,omp_target} = 25 specs (from 135 total); KF=2
% src: XSBench: 4; RSBench: 4; mixbench: 3 (all verified via ls specs/{suite}-*.json)
% src: Sums: 22+10+1+1+1=35 kernels, 60+25+4+4+3=96 specs, 53+23+4+4+3=87 PASS, 7+2+0+0+0=9 KF
\begin{tabular}{@{}lccccl@{}}
\toprule
Suite & Kernels & Total & PASS & KF & APIs \\
\midrule
Rodinia   & 22 & 60 & 53 & 7 & CUDA, OMP, OCL \\
HeCBench (curated) & 10 & 25 & 23 & 2 & CUDA, OMP, OMP-target \\
XSBench   & 1  & 4  & 4  & 0 & CUDA, OMP, OCL, OMP-target \\
RSBench   & 1  & 4  & 4  & 0 & CUDA, OMP, OCL, OMP-target \\
mixbench  & 1  & 3  & 3  & 0 & CUDA, OMP, OCL \\
\midrule
\textbf{Total} & \textbf{35} & \textbf{96} & \textbf{87} & \textbf{9} & \\
\bottomrule
\end{tabular}
\end{table*}

% src: 96=60+25+4+4+3, 87=53+23+4+4+3, 9=7+2+0+0+0 (harness/constants.py EXCLUDED_SPECS); verified 2026-04-26
% src: Rodinia 60 specs, 22 unique kernels (ls specs/rodinia-*.json); verified 2026-04-03
% src: HeCBench curated 25 specs, 10 kernels; verified 2026-04-03
% src: XSBench 4, RSBench 4, mixbench 3 specs; verified 2026-04-03
% src: 9 KNOWN_FAIL per harness/constants.py: 7 Rodinia + 2 HeCBench omp_target
The five suites contribute 96 specifications, of which 87 achieve PASS and 9 are \knownfail{} (Table~\ref{tab:suite-summary}). The 87 verified-PASS specs constitute the evaluation corpus, spanning popular computational domains including graph traversal, machine learning, stencil computation, and signal processing.

% src: specs/ on disk: rodinia=60, hecbench=25(curated), xsbench=4, rsbench=4, mixbench=3 = 96 total
% src: paper_data_together-qwen-3.5-397b-a17b.json > passk_campaign: 31 eval kernels; corpus = 35 kernels (incl. 4 KNOWN_FAIL-only)
% The evaluation corpus spans five suites: Rodinia (22 kernels, 60 specs),
% XSBench (1 kernel, 4 specs), RSBench (1 kernel, 4 specs), mixbench
% (1 kernel, 3 specs), and HeCBench (10 kernels, 25 specs), totaling
% 96 specs across 35 unique kernels
% (Appendix Table~\ref{tab:suite-summary}).

%(bfs, floydwarshall, page-rank), physics simulation (hotspot, hotspot3d, cfd, srad, iso2dfd), machine learning (backprop, nn), linear algebra (lud, nw), molecular dynamics (lavamd, md), nuclear physics (xsbench, rsbench), stencil computation (stencil1d, heat2d, jacobi), signal processing (convolution1d, scan), biophysical simulation (myocyte, heartwall), particle methods (particlefilter, streamcluster), dynamic programming (pathfinder), combinatorial search (nqueen), and GPU micro-benchmarking (mixbench).

% src: paper_data_together-qwen-3.5-397b-a17b.json > passk_campaign > overall > total=426; L0 pairs = 426/3 = 142
%At L0 (unaugmented), the corpus yields 142 unique translation pairs per model across six standard translation directions, with two additional OMP-target directions evaluated as case studies.

% src: SLoC from results/analysis/sloc_analysis.json (35 kernels); verified 2026-04-03
% src: min=80 (stencil1d), max=3304 (myocyte), median=271, mean=598.4; verified 2026-04-03
% src: 31/35 (88.6%~89%) above ParEval-Repo 133 SLoC threshold; verified 2026-04-03
% src: multi-file 24/96 (25.0%); CUDA 18/35 (51%); verified 2026-04-03

Kernel complexity spans more than an order of magnitude: from 80~SLoC (stencil1d) to 3{,}304~SLoC (myocyte), with a corpus median of 271. All 35~kernels exceed the single-function granularity typical of code-generation benchmarks: 31 of 35 (89\%) surpass ParEval-Repo's 133~SLoC threshold~\cite{ParEvalRepo2025}. Multi-file translations account for 25\% of specs (24/96), testing whether models can coordinate changes across file boundaries. A full quantitative characterization by suite is provided in Appendix Table~\ref{tab:benchmark-characterization}.

% src: 38 OMP specs = 26 omp + 12 omp_target (ls specs/*-omp*.json in curated suites); verified 2026-04-03
% src: __syncthreads in 18/35 CUDA prompt_payload files (grep __syncthreads via repo_root+source_path); verified 2026-04-03
% src: 23 OpenCL specs (ls specs/*-opencl.json); 12 omp_target specs (ls specs/*-omp_target.json); verified 2026-04-03
The corpus exercises API features spanning multiple specification versions: OpenMP v1.0 basic parallel regions (28 of 38~OMP specs) through v4.5 target offload, CUDA featurs from basic synchronization primitives (\texttt{\_\_syncthreads} through Unified Memory (1 spec) and warp shuffle intrinsics (1 spec). Three specs (covering two kernels) retain deprecated texture references, contributing to their KNOWN\_FAIL status. All 23~OpenCL~v1.x command queue API. A detailed API feature inventory is provided in Appendix~\ref{sec:appendix-b}.

With the evaluation corpus established, we now describe the experimental setup for the primary evaluation campaign.

% \subsection{Selection Funnel}
% \label{sec:selection-funnel}

% The selection is principled rather than exhaustive: the spec schema permits straightforward extension to additional suites, kernels, and APIs without modification to the evaluation harness. Each kernel in the corpus was independently verified through the complete build/run/verify pipeline prior to inclusion in any LLM evaluation experiment. 


{\color{blue}
\section{Benchmark Curation}
\label{sec:benchmark-curation}

\subsection{Suite and Kernel Selection}
\label{sec:suite-selection}

We surveyed 35 open-source HPC benchmark repositories spanning suites, mini-applications, proxy applications, libraries, and microbenchmarks (Appendix~\ref{sec:appendix-b1}). A central finding is that \textbf{repository-level counting dramatically understates the available material for parallel translation evaluation}: although only 6~repositories contain both CUDA and OpenMP implementations, kernel-level analysis reveals 472~independent CUDA--OpenMP translation pairs---a 79$\times$ multiplier (Appendix Figure~\ref{fig:repo-vs-kernel}). This gap motivates ParBench's kernel-centric evaluation design.

Five suites were selected because they provide multi-API kernel equivalence, build-run-verify automation, self-checking correctness, and domain diversity; detailed criteria and exclusion rationale are provided in Appendix~\ref{sec:appendix-b2}--\ref{sec:appendix-b6}. \textbf{Rodinia}~\cite{Rodinia2009} is the largest contributor, providing 60~specs across 22~kernels over CUDA, OpenMP, and OpenCL; 53 pass baseline verification and 7 are \knownfail{}. \textbf{HeCBench}~\cite{HeCBench2023} provides the largest initial pool: from 516~kernels, a structured funnel yields 10 curated kernels with CUDA, OpenMP, and OMP-target variants, totaling 25~specs. \textbf{XSBench}~\cite{XSBench2014} and \textbf{RSBench}~\cite{RSBench2014} contribute nuclear-physics proxy kernels, each with 4~specs, while \textbf{mixbench}~\cite{mixbench2017} contributes 3~GPU roofline micro-benchmark specs. All XSBench, RSBench, and mixbench specs pass baseline verification.

\subsection{Evaluation Corpus}
\label{sec:eval-corpus}

\begin{table*}[!htbp]
\centering
\caption{Evaluation corpus. SLoC is measured on the CUDA source variant when available. Multi-File denotes specs requiring more than one translated source file. KF = \knownfail{}.}
\label{tab:suite-summary}
\small
\resizebox{\textwidth}{!}{%
\begin{tabular}{@{}lcccclcrc@{}}
\toprule
Suite & Kernels & Specs & PASS & KF & APIs & SLoC Range & Med. & Multi-File \\
\midrule
Rodinia   & 22 & 60 & 53 & 7 & CUDA, OMP, OCL & 195--3{,}304 & 334 & 21/60 (35\%) \\
HeCBench  & 10 & 25 & 23 & 2 & CUDA, OMP, OMP-tgt & 80--235 & 180 & 0/25 (0\%) \\
XSBench   & 1  & 4  & 4  & 0 & CUDA, OMP, OCL, OMP-tgt & 1{,}390 & -- & 2/4 (50\%) \\
RSBench   & 1  & 4  & 4  & 0 & CUDA, OMP, OCL, OMP-tgt & 1{,}016 & -- & 1/4 (25\%) \\
mixbench  & 1  & 3  & 3  & 0 & CUDA, OMP, OCL & 312 & -- & 0/3 (0\%) \\
\midrule
\textbf{Total} & \textbf{35} & \textbf{96} & \textbf{87} & \textbf{9} & & \textbf{80--3{,}304} & \textbf{271} & \textbf{24/96 (25\%)} \\
\bottomrule
\end{tabular}%
}
\end{table*}

The 87 verified-\pass{} specs constitute the evaluation corpus (Table~\ref{tab:suite-summary}). Kernel complexity spans more than an order of magnitude, from 80 to 3{,}304~SLoC with a median of 271. Overall, 31 of 35 kernels (89\%) exceed 133~SLoC, the scale beyond which prior repository-level translation evaluation reports 0\% pass@$k$~\cite{ParEvalRepo2025}. Multi-file translations account for 25\% of specs, with CUDA exhibiting the highest rate due to kernel/host code separation. These properties make the corpus substantially more complex than single-function code-generation benchmarks while preserving fixed infrastructure for controlled kernel-level evaluation.
}